\chapter{First Real Finite-Strain Inelastic Continuum Law}
\label{ch:finite-strain-damage-relation}

\section{Motivation}

Chapter~\ref{ch:continuum-local-integrator-boundary} completed the mechanical
integration of the continuum-local infrastructure at the real
\texttt{Material<>} boundary, but only through a synthetic regression law.
That was enough to validate the architecture, not enough to validate the
continuum inelastic route itself.

The next stable step is therefore to plug that infrastructure into an actual
finite-strain inelastic constitutive law.  The goal of this chapter is modest
and deliberate:
\begin{itemize}
  \item introduce a first real continuum inelastic model;
  \item keep all nonlinear-solver pieces injectable at compile time;
  \item preserve the existing small-strain and hyperelastic interfaces;
  \item expose a physically meaningful internal variable through snapshots and
        VTK export.
\end{itemize}

The chosen model is an isotropic scalar damage law built on top of a
hyperelastic backbone.  It is not the final concrete model of the library, but
it is a correct and useful first continuum inelastic law because it already
needs:
\begin{itemize}
  \item finite-strain kinematics,
  \item an external algorithmic state,
  \item a local update problem,
  \item solver injection,
  \item post-processing of an inelastic internal variable.
\end{itemize}

\section{Constitutive Structure}

Let $\mathbf{E}$ denote the Green--Lagrange strain, and let the undamaged
hyperelastic free energy be
\[
  \psi_0 = \psi_0(\mathbf{E}).
\]

The damaged stored energy is taken as
\[
  \psi(\mathbf{E}, d) = g(d)\,\psi_0(\mathbf{E}),
\]
where $d \in [0,1]$ is a scalar damage variable and $g(d)$ is a degradation
function.  In the present implementation,
\[
  g(d) = k_\text{res} + (1-k_\text{res})(1-d)^2,
\]
with a small residual stiffness $k_\text{res} > 0$ to avoid a completely
singular material tangent.

The second Piola--Kirchhoff stress is then
\[
  \mathbf{S} = g(d)\,\mathbf{S}_0(\mathbf{E}),
\]
with $\mathbf{S}_0 = \partial \psi_0 / \partial \mathbf{E}$.  Likewise, in
this first stable version the algorithmic tangent is taken as the degraded
undamaged tangent:
\[
  \mathbb{C} \approx g(d)\,\mathbb{C}_0.
\]

This is a secant-like algorithmic choice rather than a fully consistent
linearization of the damage evolution law.  That limitation is explicit and
acceptable for the first stable integration step.

\section{Driving Variable and History Variable}

The damage-driving quantity is chosen as the norm of the positive spectral part
of $\mathbf{E}$:
\[
  Y(\mathbf{E}) = \left\lVert \langle \mathbf{E} \rangle_+ \right\rVert,
\]
where $\langle \cdot \rangle_+$ retains only positive principal strains.

This choice prevents compressive strains from driving damage in the present
model, which is a reasonable first approximation for tensile-dominated damage.
The history variable is
\[
  \kappa_{n+1} = \max(\kappa_n, Y_{n+1}),
\]
and the scalar damage is updated as
\[
  d(\kappa) =
  \begin{cases}
    0, & \kappa \le \kappa_0, \\
    1 - \exp\left(-\dfrac{\kappa-\kappa_0}{\kappa_f}\right),
      & \kappa > \kappa_0.
  \end{cases}
\]

The constitutive state used by the relation is therefore
\[
  \alpha = \{\kappa, d\}.
\]

\section{Local Constitutive Problem}

The update is expressed as a local nonlinear problem through the generic
continuum-local framework developed in Chapter~\ref{ch:continuum-local-problem-interface}.
For the present rate-independent isotropic damage law the local residual is
particularly simple:
\[
  R(\kappa) = \kappa - \max(\kappa_n, Y_{n+1}).
\]

Its Jacobian is
\[
  \frac{\partial R}{\partial \kappa} = 1.
\]

This scalar problem is intentionally simple, but architecturally important:
the constitutive law is now expressed through the same local-problem route that
later finite-strain plasticity, damage-plasticity, nonlocal damage, or capped
models will use.  The nonlinear algorithm itself is not hard-coded inside the
law.

\section{Compile-Time Extension Points}

The new relation is open at compile time along the same axes established in the
previous chapters:
\begin{itemize}
  \item \textbf{hyperelastic backbone}: any
        \texttt{HyperelasticModelConcept} may be used;
  \item \textbf{damage-driving force}: the law accepts a compile-time
        \texttt{DrivingForceT};
  \item \textbf{damage evolution}: the scalar softening law is injected as
        \texttt{DamageEvolutionT};
  \item \textbf{local problem}: the local update is not hard-coded into the
        relation, but expressed through
        \texttt{damage::RateIndependentDamageLocalProblem};
  \item \textbf{nonlinear algorithm}: Newton is only one admissible local
        solver; the tests also inject a direct exact solver to verify that the
        solver is genuinely interchangeable.
\end{itemize}

This is precisely the architectural property needed later for research work:
the constitutive kernel does not depend on a specific Newton implementation,
and the local-problem contract is broad enough to host other algorithms,
including PETSc-based local solvers.

\section{Reference and Spatial Paths}

Because the law is continuum-aware, the same relation accepts both:
\begin{itemize}
  \item a reference path using $(\mathbf{E}, \mathbf{S})$;
  \item a spatial path using $(\mathbf{e}, \boldsymbol{\sigma})$.
\end{itemize}

Given the degraded reference stress
\[
  \mathbf{S} = g(d)\,\mathbf{S}_0(\mathbf{E}),
\]
the spatial Cauchy stress is obtained by the standard push-forward
\[
  \boldsymbol{\sigma}
  = \frac{1}{J}\mathbf{F}\mathbf{S}\mathbf{F}^T.
\]

The same logic is used for the spatial tangent, which is obtained by
push-forward of the undamaged material tangent and then degraded by $g(d)$.

\section{Integration at the Material Boundary}

The stable boundary now works as follows:
\begin{enumerate}
  \item \texttt{ContinuumElement} builds a
        \texttt{ConstitutiveKinematics<dim>} carrier.
  \item \texttt{Material<>} forwards that carrier to an injected
        \texttt{ContinuumLocalProblemIntegrator}.
  \item The integrator solves the local damage problem using an injected local
        solver.
  \item The resulting state is committed only through
        \texttt{Material::commit(...)}; predictive evaluation remains
        side-effect free.
  \item \texttt{InternalFieldSnapshot} exports the scalar damage variable for
        post-processing, and \texttt{VTKModelExporter} now emits:
        \texttt{qp\_damage} and \texttt{nodal\_damage}.
\end{enumerate}

\section{Tests Added}

The new regression file verifies:
\begin{itemize}
  \item damage remains inactive below the threshold and the relation reduces to
        the undamaged hyperelastic backbone;
  \item damage grows under tension and does not heal upon unloading;
  \item the spatial updated-Lagrangian path returns the degraded Cauchy stress
        and degraded pushed-forward tangent;
  \item the same local problem can be solved with Newton or with an exact
        custom solver injected at compile time;
  \item the relation integrates correctly through
        \texttt{Material<ThreeDimensionalMaterial>} and through
        \texttt{MaterialPoint<ThreeDimensionalMaterial>};
  \item the internal-field snapshot now exposes the scalar damage variable.
\end{itemize}

\section{What This Enables Next}

This chapter does not yet solve finite-strain concrete damage or fracture in
their final research form.  What it does is remove another architectural
barrier:
\begin{quote}
The continuum inelastic path now has a real finite-strain law with external
state, local solver injection, spatial/reference duality, and exportable
internal variables.
\end{quote}

That makes the next steps much more concrete:
\begin{itemize}
  \item finite-strain plasticity with vector-valued local unknowns;
  \item coupled damage-plasticity for concrete;
  \item crack-band or nonlocal regularisation of scalar damage;
  \item fracture-oriented continuum models built on the same local-problem
        contract.
\end{itemize}

\section{References}

\begin{thebibliography}{9}

\bibitem{simo1987damage}
J.~C. Simo and J.~W. Ju,
\emph{Strain- and stress-based continuum damage models. I. Formulation},
International Journal of Solids and Structures, 23(7), 821--840, 1987.

\bibitem{miehe1995}
C.~Miehe,
\emph{Discontinuous and continuous damage evolution in Ogden-type large-strain
elastic materials},
European Journal of Mechanics A/Solids, 14(4), 697--720, 1995.

\bibitem{simo1998damage}
J.~C. Simo and T.~J.~R. Hughes,
\emph{Computational Inelasticity},
Springer, 1998.

\bibitem{bonet2008damage}
J.~Bonet, A.~J. Gil, and R.~D. Wood,
\emph{Nonlinear Solid Mechanics for Finite Element Analysis: Statics},
2nd ed., Cambridge University Press, 2008.

\bibitem{desouza2008damage}
E.~A. de Souza Neto, D.~Peri\'{c}, and D.~R.~J. Owen,
\emph{Computational Methods for Plasticity: Theory and Applications},
Wiley, 2008.

\end{thebibliography}
